\documentclass[10pt,a4paper]{amsart}
\usepackage[margin=19mm]{geometry}
\usepackage{amsmath,amssymb,mathtools}
\usepackage[scr=rsfs,cal=euler]{mathalfa}
\usepackage{xfrac}
\usepackage[unicode,hidelinks,bookmarks=false]{hyperref}
\newcommand{\ie}{i.e.\ }
\newcommand{\cf}{cf.\ }
\newcommand{\resp}{respectively\ }
\newcommand{\Subsection}{\subsection}
\providecommand{\nobreakdash}{\nobreak\hbox{-}}
\setlength{\emergencystretch}{2em}
\multlinegap0pt
\usepackage{dsfont}
\usepackage[shortlabels]{enumitem}
\usepackage{amssymb}
\usepackage{esint}
\usepackage{xcolor}
\newtheorem{lemma}{Lemma}
\newcommand{\T}{\mathbb{T}}
\newcommand{\dnu}{\,\mathrm{d}\nu}
\title{Cheeger Bounds for Stable Phase Retrieval in Reproducing Kernel Hilbert Spaces}
\author{Hartmut F\"uhr, Max Getter}
\date{Source: arXiv:2512.24169v1 (CC BY 4.0)}
\begin{document}
\maketitle

\noindent\emph{Source and licence.} Cheeger Bounds for Stable Phase Retrieval in Reproducing Kernel Hilbert Spaces. Version arXiv:2512.24169v1 (CC BY 4.0), distributed by arXiv under the Creative Commons Attribution 4.0 International licence, http://creativecommons.org/licenses/by/4.0/, as recorded in the arXiv licence metadata for this exact version and shown on the abstract page https://arxiv.org/abs/2512.24169v1. Full licence notice: \texttt{licenses/arxiv-2512.24169-CC-BY-4.0.txt}.\par\smallskip

\noindent\textit{Editorial scope.} Counted unit: the article's lemma lem:local\_global (source0.tex lines 479--506). The article's complete original proof of it is delivered with it (source0.tex lines 508--569, one \texttt{proof} environment of 62 lines). No mathematics is written, repaired or re-derived: the statement and the proof are the article's own lines, in the article's own notation, and no retained line is substituted or re-flowed.  Retained uncounted as the source-local context the proof needs: lines 375--381.  Nothing was omitted from any retained span, so the package carries no omission record.  No retained line contains a citation, so the wrapper carries no bibliography and no bibliography entry.  No retained line contains a figure environment, an image-inclusion command or drawing code, so no image is delivered and the package declares no asset dependency.  Editorial formatting only, in addition: an AMS \texttt{amsart} preamble that restates the article's own declarations and macros used by the retained lines, namely the environments \texttt{lemma} on separate counters, together with the standard packages those lines need.  The retained environments print in this document as Lemma 1 in order of appearance, whereas the article prints the same items with the numbering of its own full document.  The complete native source of the article as distributed in the arXiv e-print for this exact version is bundled unchanged as \texttt{sources/191858/source0.tex}.
\begin{lemma} \label{lem:norm_Borel}
	If $\Psi = (\psi_\lambda)_{\lambda \in \Lambda}$ is weakly Borel measurable, then for every $f \in L^2(A)$ the map 
	\[
	\Lambda \to [0,\infty], \quad \lambda \mapsto \| f \ast \psi_\lambda^* \|_2
	\]
	is Borel measurable.
\end{lemma}

\begin{lemma}\label{lem:local_global}
	Assume that $\Psi$ satisfies the Calder\'on condition \textup{(C)}, and let $f \in L^2(A)$. 
	\begin{enumerate}
		\item[(a)] $\Lambda_f$ is a Borel-measurable union of $\sim_f$-equivalence classes. 
		\item[(b)] Let $U \subset \Lambda$ denote any Borel measurable union of $\sim_f$-equivalence classes. Define 
		\[
		f_U := \int_U f \ast \psi_\lambda^* \ast \psi_\lambda \dnu(\lambda). 
		\]
		Then $f_U \in L^2(A)$, with 
		\[
		f_U \ast \psi_{\lambda^\prime}^* = \begin{cases}
			f \ast \psi_{\lambda^\prime}^* & \lambda^\prime \in U \\ 0 & \lambda^\prime \not \in U
		\end{cases}.
		\]
		\item[(c)] Assume that $\Lambda = \bigcup_{j \in J} U_j$, with pairwise disjoint Borel sets $U_j$ that are unions of $\sim_f$-equivalence classes. Let $(\sigma_j)_{j \in J} \in \T^J$ be arbitrary. Then 
		\[
		g = \sum_{j \in J} \sigma_j f_{U_j} 
		\] converges unconditionally in $L^2(A)$, with 
		\begin{equation} \label{eqn:changesign}
			\forall j \in J \,\forall \lambda \in U_j\,:\,  W_\Psi g(\cdot,\lambda) = \sigma_j W_\Psi f (\cdot,\lambda).
		\end{equation}
		\item[(d)] Assume that $f$ is locally phase retrievable. Let $g \in L^2(A)$ with $|W_\Psi g| = |W_\Psi f|$. Then there exists a family $(\sigma_{\lambda})_{\lambda \in \Lambda} \in \T^{\Lambda}$ such that 
		\[
		\forall \lambda \in \Lambda \,:\, g \ast \psi_{\lambda}^* = \sigma_{\lambda} (f \ast \psi_{\lambda}^*).
		\] 
		Furthermore, the map $\lambda \mapsto \sigma_\lambda$ is constant on $\sim_f$-equivalence classes. 
	\end{enumerate}
\end{lemma}

\begin{proof}
	Proof of (a): Measurability of $\Lambda_f$ follows from Lemma~\ref{lem:norm_Borel}. For $\lambda \in \Lambda \setminus \Lambda_f$, one finds $[\lambda]_{\sim_f} = \{ \lambda \}$, hence $\Lambda \setminus \Lambda_f$ is a union of equivalence classes. The same then holds for $\Lambda_f$. 
	
	Proof of (b): Given $U$, let 
	\[
	F_U: A \times \Lambda \to \mathbb{C},\quad F_U(x,\lambda) =  (f \ast \psi_\lambda^*)(x) \cdot \mathds{1}_{U}(\lambda).
	\]
	Then $F_U \in L^2(A \times \Lambda)$ due to the Calder\'on condition \textup{(C)}, hence $W_\Psi^\ast(F_U)\in L^2(A)$. Now,
	\begin{align*}
		W_\Psi^*(F_U) = \int_{\Lambda} F_U (\cdot ,\lambda) \ast \psi_\lambda \dnu(\lambda) 
		= \int_U W_\Psi f(\cdot,\lambda) \ast \psi_\lambda \dnu(\lambda)
		= f_U,
	\end{align*} 
	where the integral converges in the weak sense.
	
	Picking any $\lambda' \in U$, we have 
	\begin{align*}
		f_U \ast \psi_{\lambda'}^* = \int_U W_\Psi f(\cdot,\lambda) \ast \psi_\lambda \ast \psi_{\lambda'}^* \dnu(\lambda)
		=   \int_U f \ast \psi_\lambda^* \ast \psi_\lambda \ast \psi_{\lambda'}^* \dnu(\lambda).
	\end{align*}
	Since $[\lambda']_{\sim_f} \subset U$, it follows that, for all $\lambda \in \Lambda \setminus U$, 
	\[
	\| f \ast \psi_\lambda^* \ast \psi_{\lambda'} \|_2 = 0.
	\] 
	Thus, using the weak-sense inversion formula for $f$, we find that
	\begin{align*}
		\int_U f \ast \psi_\lambda^* \ast \psi_\lambda \ast \psi_{\lambda'}^* \dnu(\lambda) 
		= \int_\Lambda f \ast \psi_\lambda^* \ast \psi_\lambda \ast \psi_{\lambda'}^* \dnu(\lambda)
		= f \ast \psi_{\lambda'}^*.
	\end{align*} 
	If instead $\lambda' \in \Lambda \setminus U$ and $\lambda \in U$, the assumption that $U$ is a union of $\sim_f$-equivalence classes yields 
	\[ f \ast \psi_{\lambda} \ast \psi_{\lambda'}^* = 0 \]
	and therefore 
	\begin{align*}
		f_U \ast \psi_{\lambda'}^* = \int_U f \ast \psi_\lambda^*  \ast \psi_\lambda \ast \psi_{\lambda'}^* \dnu(\lambda)  = 0. 
	\end{align*}
	This shows part (b). 
	
	For part (c), part (b) implies that $(\sigma_j f_{U_j})_{j \in J}$ is a family in $L^2(A)$. Pairing part (b) also with the isometry property of $W_\Psi$ yields, for all $j \not= k$,  
	\[
	\langle f_{U_j}, f_{U_k} \rangle = \langle W_\Psi f_{U_j}, W_\Psi f_{U_k} \rangle = 0. 
	\] 
	Since the $U_j$ are disjoint and cover all of $\Lambda_f$, we obtain 
	\[
	\sum_{j \in J} \|f_{U_j} \|_2^2 = \sum_{j \in J} \| W_\Psi f_{U_j} \|_2^2 = \| W_\Psi f \|_2^2 = \| f \|_2^2.
	\]
	Hence
	\[
	g = \sum_{j \in J} \sigma_j f_{U_j}
	\] is an unconditionally converging orthogonal sum, for every choice $(\sigma_j)_{j \in J} \in \T^J$, and (\ref{eqn:changesign}) follows from (b) and disjointness of the $U_j$.
	
	
	Proof of (d): 
	The existence of the family $(\sigma_\lambda)_{\lambda \in \Lambda}$ with the desired properties follows from the definition of local phase retrievability. 
	
	For $\lambda \not= \lambda'$, associativity and commutativity of convolution allow to compute
	\begin{align*}
		f \ast \psi_\lambda^\ast \ast \psi_{\lambda'}^\ast = \sigma_\lambda^{-1} (g \ast \psi_\lambda^\ast \ast \psi_{\lambda'}^\ast) 
		= \sigma_{\lambda'} \sigma_{\lambda}^{-1} (f \ast \psi_\lambda^\ast \ast \psi_{\lambda'}^\ast),
	\end{align*}
	which for $\lambda \approx_f \lambda'$ entails $\sigma_\lambda = \sigma_{\lambda'}$. Since $\sim_f$ is the transitive hull of $\approx_f$, we get $\sigma_\lambda = \sigma_{\lambda'}$ for all $\lambda' \in [\lambda]_{\sim_f}$. 
\end{proof}

\end{document}
